\section{Conclusion}
\label{sec:conclusion}

This paper develops a Copula-based theory of mean-reverting Markov processes
and studies its empirical content without equating stationarity, dependence
decay, and a restoring force.  The corrected framework separates a
drift-reverting transition kernel from mean convergence under an initial law,
requires strong conditional-mean persistence to be contractive, and fixes the
almost-everywhere versions needed for Copula densities and conditional kernels.
Strict monotone diffusion transforms yield observable densities only under
invertibility and support conditions.

The stationary-reset family is the central construction.  Within the scalar
mixture class, Chapman--Kolmogorov consistency forces multiplicative weights
whenever the base kernel is not already the invariant projection.  A regular
cumulative hazard gives the familiar exponential survival form; an ordinary
intensity requires absolute continuity, not merely right continuity.  On the
centered $L^2$ space, reset mixing multiplies the base semigroup and every
stationary covariance by $e^{-\lambda t}$.  Under a reversible base semigroup,
the spectral gap shifts by exactly $\lambda$.  These results clarify both the
advantage and the identification limit of the construction: conditional-mean
and second-order decay see $\kappa+\lambda$, while the separate rates must come
from transition shape or coherent multi-horizon restrictions.

The empirical revision changes the substantive conclusion.  Exact ties are a
first-order feature of Kraken stablecoin prices, not a minor numerical issue.
They make the continuous mixed-OU point-density likelihood unbounded and
invalidate the earlier continuous-path LR bootstrap.  A nearest-tick interval
likelihood removes that divergence.  Using separated full and frozen-training
estimates, it finds positive out-of-sample MOU--OU score differences for USDT
and USDC and weaker evidence for DAI.  The common calendar also reveals marked
cross-asset heterogeneity, but those point estimates are not yet formal
mechanism comparisons.

Most importantly, one-minute dynamics do not transfer successfully to the
five-minute kernel for any of the three assets.  The continuous-time semigroup
restriction therefore fails under the present interval approximation.  The
evidence supports dependence refresh as a useful predictive representation;
it does not yet validate frequency-invariant structural rates or identify
arbitrage, redemption, or collateral mechanisms.  The next decisive step is
not a larger continuous bootstrap.  It is scalable full filtering of the
rounded observation model, followed by multi-horizon estimation and a
synchronous day-block cross-asset bootstrap.  Structural economic
interpretation should remain conditional on those gates.
